% Standalone source-record documentation; no research attempts.
\documentclass[11pt]{article}
\usepackage[margin=25mm]{geometry}
\usepackage{fontspec,unicode-math}
\setmainfont{Latin Modern Roman}
\setmathfont{Latin Modern Math}
\usepackage{amsmath,amssymb,mathtools,xurl,hyperref}
\hypersetup{colorlinks=true,urlcolor=blue,linkcolor=blue}
\setcounter{secnumdepth}{0}
\setlength{\parindent}{0pt}
\setlength{\parskip}{5pt}
\setlength{\emergencystretch}{3em}
\newcommand{\sref}[2]{\hyperlink{source:#1:#2}{[S#2]}}
\newcommand{\cataloguescope}[1]{\paragraph{Source record.}#1}
\newcommand{\floor}[1]{\left\lfloor#1\right\rfloor}
\begin{document}
% UnsolvedMath ID 30000196; source catalogue code OWR-792-010
\setcounter{section}{30000195}
\section{Observability for Polynomial Systems}\label{30000196}
{\small\sffamily UnsolvedMath ID 30000196 · Computer Science}
\subsection{Definitions and mathematical statement}

\paragraph{Shared notation.}Standard number systems are $\mathbb N,\mathbb Z,\mathbb Q,\mathbb R,\mathbb C$. Specialized notation and index ranges follow the cited source and the definitions below.


A polynomial system has state and output dynamics given by polynomial maps. Observability concerns distinguishing initial states from outputs; algebraic observability uses algebraic information in the output and its derivatives.
\paragraph{Record scope and qualifications.}
The catalogue title is retained as metadata. Named certificates, contracts, bounds or frameworks in that title are not asserted as results of the original source extract.
The family of systems, coefficient topology, and precise meaning of generic must be supplied by the report.

Source-record limitation: The genericity family and algebraic-observability criterion are absent, while the definition-choice clause is a broad agenda.
\paragraph{Statement recorded in the supplied source.}
What is the right definition of observability for polynomial systems? Moreover, is the more restricted condition of algebraic observability generic? \sref{30000196}{2}
\paragraph{Source status.}
The supplied catalogue labels this source record partially solved; that metadata was checked against the supplied snapshot on 2026-10-11. The label is preserved for provenance and is not an independent verification of present open status or mathematical correctness. The catalogue literature assessment is dated 2026-08-21.
\subsection{Short English statement}
Identify an appropriate observability concept for polynomial systems and decide whether algebraic observability is generic.
\subsection{Sources}
\begin{itemize}
\item[\hypertarget{source:30000196:1}{S1}] Ulam AI and the UnsolvedMath Contributors, UnsolvedMath: A Curated Collection of Open Mathematics Problems, record 30000196 (OWR-792-010). Catalogue curation and metadata: CC BY 4.0. The source extract is retained; definitions, English summaries and qualifications are editorial additions.
\url{https://huggingface.co/datasets/ulamai/UnsolvedMath}.
\item[\hypertarget{source:30000196:2}{S2}] Regelungstheorie (2005), Oberwolfach Reports, DOI 10.4171/owr/2005/11, pp. 565--645
\url{https://doi.org/10.4171/owr/2005/11}.
\item[\hypertarget{source:30000196:3}{S3}] generic algebraic-observability result
\url{https://www.researchgate.net/publication/228999815_Genericity_of_Algebraically_Observable_Polynomial_Systems}.
\item Source-specific Oberwolfach Reports licences, including CC BY-SA 4.0 where applicable, are retained; see THIRD\_PARTY\_NOTICES.md.
\end{itemize}
\paragraph{Preparation and review.}
The original source and UnsolvedMath contributors retain attribution for the source extract. Codex (GPT-6; exact serving revision not exposed) assisted with definitions, English summaries, scope notes and formatting. No new human mathematical authorship or independently verified solution is asserted. Permanent source ID 30000196 and code OWR-792-010 checked against the complete supplied registry and current repository. Distinct upstream records, including repeated questions and broad or damaged extracts, are retained for complete coverage of the source CS category. A repeated record is not a new independent conjecture. Complete mathematical-equivalence and current-literature review has not been performed. Preparation and final review follow contribution prompts C and D, with source gaps disclosed. Structural and rendering checks do not establish a new theorem or open status.
\label{end:30000196}
\end{document}
